% ---------------------------------------------------------------------
%  PHẦN MỞ ĐẦU (không đánh số chương)
% ---------------------------------------------------------------------
\chapter*{Mở đầu}
\phantomsection
\addcontentsline{toc}{chapter}{Mở đầu}
\markboth{Mở đầu}{}
\label{ch:modau}

\begin{muctieu}
Chương mở đầu trình bày bối cảnh và lý do lựa chọn đề tài xây dựng hệ thống
hỗ trợ sàng lọc và phân tầng nguy cơ y tế, xác lập các mục tiêu kỹ thuật cụ
thể, giới hạn phạm vi bài toán trong khuôn khổ phân tầng nguy cơ và tóm lược
bố cục báo cáo.
\end{muctieu}

\section*{Lý do chọn đề tài}
\phantomsection
\addcontentsline{toc}{section}{Lý do chọn đề tài}

Trong bối cảnh hệ thống y tế cơ sở tại Việt Nam thường xuyên đối mặt với tình
trạng quá tải bệnh nhân, nhu cầu tự tìm kiếm thông tin y tế và tự sử dụng thuốc
không kê đơn \en{(Over-the-counter - OTC)} của người dân ngày càng gia tăng.
Theo Tổ chức Y tế Thế giới \cite{WHO2021}, việc thiếu định hướng y tế ban đầu
dẫn tới hai hệ quả tiêu cực phổ biến: hoặc người bệnh chủ quan bỏ sót các
triệu chứng cảnh báo nguy kịch \en{(red flags)}, dẫn đến chậm trễ thời gian
vàng cấp cứu; hoặc người bệnh hoang mang, đi khám vượt tuyến không cần thiết,
gây lãng phí nguồn lực xã hội.

Sự bùng nổ của các mô hình ngôn ngữ lớn \en{(Large Language Models - LLMs)}
gần đây mở ra tiềm năng to lớn trong việc giao tiếp tự nhiên và giải thích
thông tin sức khỏe \cite{Singhal2023, Thirunavukarasu2023}. Tuy nhiên, việc áp
dụng trực tiếp các mô hình AI tạo sinh đơn thuần vào bài toán ra quyết định y tế
tiềm ẩn nguy cơ nghiêm trọng về mặt an toàn người bệnh:
\begin{itemize}[leftmargin=1.6em, itemsep=2pt]
  \item \term{Hiện tượng ảo giác \en{(Hallucination)}:} Mô hình có thể tự bịa đặt
        liều lượng, tương tác thuốc hoặc bỏ qua các triệu chứng cảnh báo cấp
        cứu, sinh ra các tư vấn sai lệch nhưng với giọng điệu tự tin.
  \item \term{Tính thiếu nhất quán \en{(Non-determinism)}:} Với cùng một tập triệu
        chứng lâm sàng, mô hình tạo sinh có thể đưa ra các mức độ xử trí khác
        nhau tại các lần chạy, vi phạm nguyên tắc an toàn y khoa.
  \item \term{Thiếu kiểm soát hộp đen \en{(Black-box nature)}:} Các bác sĩ và dược
        sĩ chuyên môn không thể truy vết hoặc kiểm soát tuyệt đối biên giới an
        toàn của các mô hình mạng nơ-ron sâu \cite{Musen2014}.
\end{itemize}

Để giải quyết bài toán này một cách khoa học và an toàn, đề tài lựa chọn phương
pháp tiếp cận \term{Kiến trúc lai \en{(Hybrid AI Architecture)}}: kết hợp sức mạnh
của \textbf{Hệ chuyên gia dựa trên tri thức luật tất định} \en{(Deterministic
Rule-based System)} \cite{Buchanan1984} với khả năng tương tác linh hoạt, trích
xuất sự kiện và diễn giải người dùng của \textbf{Mô hình ngôn ngữ lớn có cấu trúc
hóa} \en{(Function/Tool Calling)} \cite{Schick2023}. Cách tiếp cận này đảm bảo
rằng mọi quyết định sinh tử (cấp cứu, chuyển tuyến) đều được neo chặt vào các quy
tắc lâm sàng đã được kiểm chứng, trong khi người bệnh vẫn có được trải nghiệm tư
vấn tự nhiên, ân cần và dễ hiểu.

\section*{Mục tiêu của đồ án}
\phantomsection
\addcontentsline{toc}{section}{Mục tiêu của đồ án}

Đề tài đặt ra các mục tiêu cụ thể, đo lường được nhằm hoàn thiện giải pháp:
\begin{itemize}[leftmargin=1.6em, itemsep=2pt]
  \item \textbf{Mục tiêu 1 — Chuẩn hóa cơ sở tri thức phân tầng y tế:} Xây dựng
        tập dữ liệu lâm sàng gồm tối thiểu 30 triệu chứng, 20 dấu hiệu cờ đỏ cấp
        cứu và 40 quy tắc phân tầng nguy cơ theo tiêu chuẩn phân tầng cấp cứu
        quốc tế \cite{Mackway2014}, hỗ trợ nhận diện 4 mức độ: Cấp cứu
        \en{(Emergency)}, Đi khám bác sĩ \en{(Refer Doctor)}, Tư vấn dược sĩ
        \en{(Pharmacist Review)} và Tự chăm sóc tại nhà \en{(Self-care)}.
  \item \textbf{Mục tiêu 2 — Thiết kế và cài đặt động cơ phân tầng an toàn:} Xây
        dựng hệ thống ra quyết định theo nguyên tắc bảo thủ lâm sàng
        \en{(Conservative Evaluation)}, đảm bảo trên mọi kịch bản kiểm thử không
        có ca cấp cứu nào bị xếp xuống mức xử trí nhẹ hơn, cơ chế phân tầng có
        tính tất định và không bị hạ cấp nguy cơ bởi kết quả trích xuất của AI.
  \item \textbf{Mục tiêu 3 — Xây dựng mô-đun diễn giải lâm sàng bằng Tool-calling:}
        Hiện thực hóa dịch vụ sinh lời giải thích \en{(AI Briefing)} qua schema
        gọi hàm có cấu trúc, tích hợp bộ kiểm duyệt chặt chẽ \en{(Sanitization)}
        loại bỏ nội dung mỏng, đảm bảo luôn có lý do lâm sàng và tối thiểu hai
        danh mục hướng dẫn sức khỏe thiết thực.
  \item \textbf{Mục tiêu 4 — Hiện thực hóa hệ sinh thái tương tác thực tế:} Triển
        khai giải pháp thành ứng dụng web hoàn chỉnh trên nền tảng Laravel, hỗ
        trợ xử lý bất đồng bộ chống nghẽn \en{(Async Jobs)}, cơ chế phục hồi
        đồng bộ \en{(Fallback Seam)} và hàng đợi cho dược sĩ duyệt ca.
\end{itemize}

\section*{Phạm vi và đối tượng}
\phantomsection
\addcontentsline{toc}{section}{Phạm vi và đối tượng}

Phạm vi thực hiện của đồ án được xác định rõ ràng như sau:
\begin{itemize}[leftmargin=1.6em, itemsep=2pt]
  \item \textbf{Đối tượng phục vụ:} Người dùng đại chúng có các triệu chứng bệnh
        thông thường (đau đầu, sốt, tiêu hóa, hô hấp, cơ xương khớp nhẹ) tìm
        kiếm hướng dẫn chăm sóc sức khỏe ban đầu.
  \item \textbf{Phạm vi chuyên môn:} Hệ thống chỉ thực hiện \term{sàng lọc và phân
        tầng nguy cơ} \en{(Triage \& Risk Stratification)}, không thực hiện chẩn
        đoán bệnh danh y khoa chuyên sâu và không kê đơn thuốc đặc trị.
  \item \textbf{Ngôn ngữ và dữ liệu:} Hỗ trợ ngôn ngữ tiếng Việt tự nhiên; danh
        mục thuốc gợi ý chỉ giới hạn trong nhóm thuốc không kê đơn \en{(OTC)}
        hiện diện trong cơ sở dữ liệu nhà thuốc và đã đối chiếu hồ sơ chống chỉ
        định của người bệnh.
  \item \textbf{Giới hạn kỹ thuật:} Hệ thống hoạt động dưới hình thức tương tác
        văn bản trên giao diện web, dữ liệu lưu trữ tại cơ sở dữ liệu quan hệ nội
        bộ, kết nối mô hình ngôn ngữ lớn thông qua giao thức API chuẩn hóa.
\end{itemize}

\section*{Bố cục báo cáo}
\phantomsection
\addcontentsline{toc}{section}{Bố cục báo cáo}

Nội dung báo cáo được kết cấu thành phần Mở đầu và 5 chương chính:
\begin{itemize}[leftmargin=1.6em, itemsep=0pt]
  \item \textbf{Phần Mở đầu:} Đặt vấn đề, lý do nghiên cứu, mục tiêu, phạm vi và cấu trúc báo cáo.
  \item \textbf{Chương~\ref{ch:coso} — Cơ sở lý thuyết:} Tổng quan về hệ chuyên gia
        y tế, phân tầng lâm sàng, mô hình ngôn ngữ lớn và kỹ thuật gọi hàm.
  \item \textbf{Chương~\ref{ch:thietke} — Phân tích và thiết kế:} Mô hình hóa yêu
        cầu, thiết kế kiến trúc hệ thống Hybrid AI Pipeline, thiết kế cơ sở tri
        thức y khoa và schema kiểm soát an toàn.
  \item \textbf{Chương~\ref{ch:caidat} — Cài đặt:} Hiện thực hóa các dịch vụ lõi,
        quy trình xử lý công việc bất đồng bộ, bộ kiểm duyệt dữ liệu và cơ chế
        cứu hộ tự động.
  \item \textbf{Chương~\ref{ch:ketqua} — Kết quả và đánh giá:} Báo cáo chỉ số
        thực nghiệm, khả năng phát hiện cờ đỏ và thời gian đáp ứng.
  \item \textbf{Chương~\ref{ch:ketluan} — Kết luận:} Đánh giá mức độ hoàn thành so
        với mục tiêu, các bài học rút ra và hướng phát triển tiếp theo.
\end{itemize}
